\subsection{Equation of a Circle Given a Diameter} 
\begin{definition}[Diameter]
A \term{diameter} of a circle is a line segment passing through the center of the circle with end points on the edge. \emph{The} diameter of a circle is the length of such a segment.
\begin{icpic}
\fill (0,0) circle (0.1);
\draw[geom] (30:-2) -- (30:2);
\draw[geom,blue] (0,0) circle (2);
\foreach \r in {-2,2} \fill (30:\r) circle (0.1);
\end{icpic}
\end{definition}
\begin{note}
The segment on each side of the center is a \makeblank[1in], so each must be the same length. So the center is the \makeblank[1in] of any diameter.
\end{note}
\begin{note}
The diameter $d$ of a circle is \makeblank[1in] its radius $r$. In other words, the radius is \makeblank[1in] the diameter.
\end{note}
We can write the equation of a circle given the endpoints of a diameter.
\begin{example}
Give the equation of a circle whose diameter has endpoints $A=(-3,1)$ and $B=(5,2)$.
\begin{itemize}
\item The center is at the midpoint of \makeblanksmall.
\item $h=\makeblank[1in]=\makeblanksmall$
\item $k=\makeblank[1in]=\makeblanksmall$
\item The diameter is the length of \makeblanksmall.
\item $d=\makeblank=\makeblanksmall$
\item $r=\frac{d}{2}$, so $r^2=\makeblanksmall=\makeblanksmall$
\item In this example, $r^2=\makeblanksmall$
\item The equation is \makeblank.
\end{itemize}
\end{example}